% ===================================================================
%  اللوحة رقم ١ — المدرسة. — الثانوية. — الفصل.
%
%  Ceci n'est PAS une transcription : c'est une traduction, faite en
%  2026, en arabe d'aujourd'hui. D'ou trois differences avec texto/io
%  et texto/fr, et trois seulement :
%
%   * pas de \nl ni de \cc — la traduction ne suit aucune ligne
%     imprimee, et les lignes de ce fichier ne sont que du confort ;
%   * pas de \begin{VUpage} — elle n'a ni feuillet ni folio, et la
%     page de lecture ne lui en affiche pas ;
%   * le gras se pose sur le TERME seul, ponctuation dehors.
%
%  Tout le reste est identique, et doit l'etre : les cles « %%K » sont
%  celles de texto/io, macro pour macro pour l'apparat, et les renvois
%  \textsuperscript{(n)} visent les MEMES objets de la planche.
%
%  L'ido est le texte de base ; le francais sert de controle, et
%  l'anglais et l'espagnol de calibrage.
%
%  UNE LANGUE QUI S'ECRIT DE DROITE A GAUCHE. Le fichier, lui, se lit
%  dans l'ordre logique — le premier caractere d'un bloc est le
%  premier mot de la phrase — et c'est le navigateur qui retourne
%  l'affichage : html.py pose « dir="rtl" » sur la seule case arabe.
%  Les renvois restent donc ecrits \textsuperscript{(18)}, avec leurs
%  chiffres latins, APRES le substantif comme dans les six autres
%  colonnes ; l'algorithme bidirectionnel les remet du bon cote tout
%  seul. Ne rien inverser a la main : ce qui est retourne deux fois
%  revient a l'endroit.
%
%  L'ordinal se place APRES le nom — « السلسلة الأولى », « المشهد
%  الأول » — d'ou un membre a part dans les motifs SERIO et CENO de
%  html.py, la ou les langues latines n'ont demande qu'un mot de plus.
%
%  LE NOM DE L'EDITEUR SE PRONONCE « del-MASS ». Delmas est un nom du
%  Midi -- « del mas », celui du mas -- et son s final se fait
%  entendre, comme dans Chaban-Delmas ; l'imprimerie est d'ailleurs a
%  Bordeaux, rue Saint-Christoly. La transcription دِلما le disait
%  « Delma » et perdait la syllabe : c'est دِلماس.
% ===================================================================

%%K t01-apar-0 apar
\VUsaut{2.0mm}
\VUcentre{12.6pt}{120}{كُتيِّب الشرح}
\VUsaut{3.2mm}
\VUcentre{8.4pt}{160}{الملحق بـ}
\VUsaut{2.6mm}
\VUtitre{15.6pt}{0}{لوحات دِلماس المساعِدة}
\VUsaut{3.4mm}
\VUornamento{9pt}
\VUsaut{5.0mm}
\VUcentre{13.2pt}{90}{السلسلة الأولى}
\VUsaut{3.0mm}
\VUfilet{16mm}
\VUsaut{5.6mm}

%%K t01-apar-1 apar
\VUtitre{15.0pt}{60}{اللوحة رقم ١}
\VUsaut{1.4mm}
\VUtitre{11.4pt}{0}{المدرسة. --- الثانوية. --- الفصل.}
\VUsaut{2.2mm}
\VUfilet{20mm}
\VUsaut{6.4mm}

%%K t01-tit-1 sub
\VUpk{10.2pt}{40}{وصف عام.}
\VUsaut{3.0mm}

%%K t01-01-1 p
1. --- تُظهِر هذه اللوحة \VUgras{فصلًا} في \VUgras{ثانوية}. في هذا الفصل
ثمانية \VUgras{طاولات} \textsuperscript{(18)} وثمانية \VUgras{مقاعد}
\textsuperscript{(32)}. على كل مقعد يجلس ثلاثة تلاميذ. ويجلس \VUgras{المعلم}
\textsuperscript{(7)} على \VUgras{كرسي} \textsuperscript{(8)} أمام \VUgras{منبره}
\textsuperscript{(6)}.

%%K t01-02-1 p
2. --- على يسار المنبر تقوم \VUgras{خزانة} \textsuperscript{(15)} ذات واجهة
زجاجية. وعلى اليمين \VUgras{السبورة} \textsuperscript{(1)} ومعها \VUgras{الممسحة}
\textsuperscript{(2)} و\VUgras{الطباشير} \textsuperscript{(3)}.

%%K t01-02-2 p
وفوق المنبر تُعلَّق \VUgras{لوحات جدارية} \textsuperscript{(9, 11, 12)}
و\VUgras{خريطة} \textsuperscript{(10)}.

%%K t01-03-1 p
3. --- ليس كل التلاميذ في \VUgras{أماكنهم} \textsuperscript{(17)}. يوحنا
\textsuperscript{(4)} واقف عند السبورة؛ يمسك الممسحة ويمحو \VUgras{الأشكال
الهندسية}.

%%K t01-03-2 p
وبطرس قرب المنبر؛ يجمع \VUgras{الواجبات المكتوبة} \textsuperscript{(5)} ويحملها
إلى المعلم.

%%K t01-04-1 p
4. --- \VUgras{هذا} المعلم هو معلم \VUgras{اللغات الحديثة}. يعلّم التلاميذ
أن يتكلموا اللغات الأجنبية ويقرأوها ويكتبوها \VUgras{(الألمانية
والإنجليزية والإسبانية والإيطالية والروسية} أو \VUgras{الفرنسية)}.

%%K t01-05-1 p
5. --- الفصل واسع جدًا ومضيء جدًا. في آخره \VUgras{نافذة} عريضة لها
\VUgras{عارضتان} \textsuperscript{(78)} و\VUgras{باب زجاجي}، لتهويته وإنارته.

%%K t01-05-2 p
وهذا الفصل ليس باردًا. ففي وسطه، لتدفئته، \VUgras{مِدفأة} \textsuperscript{(46)}
جيدة جدًا مع \VUgras{مِدخنتها} \textsuperscript{(45)}.

%%K t01-05-3 p
وفي الحاجز \VUgras{مَشاجب} \textsuperscript{(58)} مثبَّتة؛ تُعلَّق فيها قبعات
التلاميذ ومعاطفهم.

%%K t01-tit-2 sub
\VUsaut{5.2mm}
\VUpk{10.2pt}{40}{الفناء.}
\VUsaut{3.6mm}

%%K t01-06-1 p
6. --- \VUgras{الساعة} \textsuperscript{(63)} تشير إلى التاسعة وخمس دقائق.

%%K t01-06-2 p
انتهت \VUgras{الفسحة} \textsuperscript{(76)} للتو وعاد الدرس يبدأ. والفسحة تكون
في \VUgras{الفناء} \textsuperscript{(75)}. ونرى بعض التلاميذ يلعبون. ويراقبهم
\VUgras{مراقب} \textsuperscript{(74)}.

%%K t01-07-1 p
7. --- وفي الفناء نرى أيضًا أشخاصًا آخرين: \VUgras{المفتش}
\textsuperscript{(73)} و\VUgras{المدير} \textsuperscript{(71)} و\VUgras{الناظر}
\textsuperscript{(72)}.

%%K t01-07-2 p
المفتش يفتش المدارس، والمدير يدير المؤسسة، والناظر يتابع عمل التلاميذ.

%%K t01-07-3 p
والرابع هو \VUgras{الحاجب} \textsuperscript{(77)}. يحمل تحت إبطه سجلًا ليقيّد
الغائبين.

%%K t01-08-1 p
8. --- وفي آخر الفناء \VUgras{أجهزة الرياضة} \textsuperscript{(70)} وورشة
\VUgras{الأشغال اليدوية} \textsuperscript{(69)}.

%%K t01-08-2 p
وعلى يمين اللوحة نلمح \VUgras{فصلي} \VUgras{الموسيقى} و\VUgras{الرسم}.

%%K t01-noto-1 noto
\VUnotes{143.2mm}{%
(*) أما \textit{الأسماء الأولى} فننصح كذلك بنقلها في صورتها اللاتينية.
فبذلك وحده تُتجنَّب صيغ لا تُحصى. ثم كيف يُختار بين \textit{Giovanni,
Jean, Johann, Jan, Hans, John, Ivan} وغيرها؟ أنضع معجمًا خاصًا للأسماء؟
لنأخذ صيغة الرفع اللاتينية ولنقل: \VUgras{Ioannes, Ioanna; Iulius,
Iulia; Iakobus; Andreas; Lukas.} (النحو الكامل).%
}

%%K t01-tit-3 sub
\VUpk{10.2pt}{40}{وصف مفصَّل.}
\VUsaut{2.4mm}

%%K t01-tit-4 sub
\VUpk{10.2pt}{40}{التلاميذ المجتهدون.}
\VUsaut{3.0mm}

%%K t01-09-1 p
9. --- التلاميذ الجالسون على المقعد الأول يمين المنبر هم \VUgras{هنري}
(*) و\VUgras{إسطفان} و\VUgras{شارل}.

%%K t01-09-2 p
هنري منتبه جدًا ومجتهد؛ يصغي إلى المعلم. وعلى طاولته \VUgras{كُرّاسة}
\textsuperscript{(28)} و\VUgras{كتاب} \textsuperscript{(29)} و\VUgras{معجم}
\textsuperscript{(30)} و\VUgras{مقلمة} \textsuperscript{(31)}. وكتابه فيه
\VUgras{صفحات} كثيرة؛ وكل فصل فيه عدة \VUgras{فقرات}، وكل \VUgras{سطر} فيه
\VUgras{كلمات} كثيرة.

%%K t01-09-4 p
وجاره إسطفان \textsuperscript{(24)} نشيط. يدوّن في كُرّاسته الدرس المطلوب للمرة
القادمة. وله \VUgras{خط} جميل. وكتابه مغلق. لا يُرى منه إلا \VUgras{الغلاف}
\textsuperscript{(27)}. وإلى جانبه \VUgras{فِرجاره} \textsuperscript{(26)}.

%%K t01-09-5 p
وشارل \textsuperscript{(23)} يمسك كتابه بيده؛ يفتحه ليراجع درسه مرة أخرى. وأمامه،
كما أمام كل تلميذ، \VUgras{مِحبرة} \textsuperscript{(25)}. وهو مطيع.

%%K t01-10-1 p
10. --- وعلى الطاولة التالية \VUgras{لويس وأنطون} و\VUgras{بولس}. لويس
يُسمِّع \VUgras{درسه} \textsuperscript{(22)} ويجيب عن \VUgras{أسئلة} المعلم.
وأنطون \textsuperscript{(21)} قليل الانتباه جدًا؛ بل يعاقبه المعلم أحيانًا. وبولس
\textsuperscript{(20)} رفيق طيب، لطيف وخدوم. وهو منتبه جدًا.

%%K t01-tit-5 sub
\VUsaut{4.2mm}
\VUpk{10.2pt}{40}{التلاميذ المهملون.}
\VUsaut{3.2mm}

%%K t01-11-1 p
11. --- خلف هنري يجلس \VUgras{جورج} \textsuperscript{(38)}. ليس ولدًا سيئًا،
لكنه \VUgras{ثرثار} وقليل الانتباه وكثير الحركة. \VUgras{طيشه}
و\VUgras{عصيانه} يُحدثان \VUgras{فوضى} في الدرس، وكثيرًا ما يستحق
\VUgras{تنبيهًا} شديدًا. \VUgras{مسوّدته} \textsuperscript{(35)} متسخة، يلطخها
\VUgras{ببقع حبر} من \VUgras{قلمه} \textsuperscript{(34)}، ولذلك يستعمل
\VUgras{المِمحاة} \textsuperscript{(36)} دائمًا؛ و\VUgras{حقيبته} \textsuperscript{(33)}
ممزقة، فهو مهمل جدًا. ولمَ يحمل \VUgras{ماسك القلم} \textsuperscript{(37)} إلى
فصل اللغات الحديثة؟

%%K t01-11-2 p
وجاره إدموند \textsuperscript{(39)} لا يحبه. صحيح أنه كسول قليلًا، لكنه هادئ
ساكن. لا يثرثر؛ غير أنه، بدل أن يصغي، يؤثر أن يقرأ \VUgras{الحكايات} في
\VUgras{كتاب القراءة}.

%%K t01-11-3 p
والتلميذ الثالث على هذا المقعد هو \VUgras{لودوفيك} \textsuperscript{(40)}. وهو
مواظب. يكتب على \VUgras{ورقة} بيضاء. وقد بسط أمامه \VUgras{ورق النشاف}
\textsuperscript{(41)}.

%%K t01-12-1 p
12. --- التلاميذ الجالسون على المقعد الثاني، يسار المعلم، صغار جدًا.
فأولهم \VUgras{إميل} الصغير لا يحسن الكتابة بعد؛ يحمل \VUgras{لوحه}
\textsuperscript{(42)} و\VUgras{قلمه الحجري} و\VUgras{إسفنجته} \textsuperscript{(43)}.

%%K t01-12-2 p
وإلى جانبه \VUgras{أوغست} و\VUgras{برنار} \textsuperscript{(44)}. وهذا الأخير
يعمل جيدًا، لكن \VUgras{هيئته} غير مرتبة أبدًا.

%%K t01-13-1 p
13. --- وخلفهم يجلس ثلاثة \VUgras{كسالى}. \VUgras{ألبير} لا يحفظ دروسه.
و\VUgras{يعقوب} \textsuperscript{(47)} نائم بدل أن يصغي. \VUgras{كسله} يجلب له
\VUgras{التوبيخ} بل \VUgras{العقاب}. وأخيرًا \VUgras{كريستيان}
\textsuperscript{(48)} طائش أكثر مما ينبغي. \VUgras{سلوكه} سيئ ولا يجتهد في
إصلاحه.

%%K t01-14-1 p
14. --- أما جيرانهم فحسنو السيرة. \VUgras{إرنست} \textsuperscript{(51)} يحب
النظام والانضباط؛ يستعمل دائمًا \VUgras{مسطرته} \textsuperscript{(50)}
و\VUgras{قلمه الرصاص} \textsuperscript{(49)}. وهو \VUgras{تلميذ خارجي}.

%%K t01-14-2 p
و\VUgras{فرنسيس} \textsuperscript{(52)} \VUgras{تلميذ داخلي}؛ يلبس الزي ويأكل
وينام في المؤسسة. وهو \VUgras{رفيق} طيب.

%%K t01-14-3 p
وموريس يبري \VUgras{قلمه الرصاص} \VUgras{بمُدْيته} \textsuperscript{(56)}؛ وهو
شديد العناية.

%%K t01-14-4 p
يحمل كتبه كلها، و\VUgras{نحوه} \textsuperscript{(55)} و\VUgras{دفتره}
\textsuperscript{(54)}، بل و\VUgras{مفكرة} \textsuperscript{(53)}. و\VUgras{محفظته}
\textsuperscript{(57)} على الأرض خلفه.

%%K t01-14-5 p
والطاولتان في آخر الفصل يشغلهما \VUgras{تلاميذ مجتهدون} \textsuperscript{(19)}.

%%K t01-tit-6 sub
\VUsaut{4.6mm}
\VUpk{10.2pt}{40}{الرسم والموسيقى.}
\VUsaut{3.4mm}

%%K t01-15-1 p
15. --- في \VUgras{فصل الرسم} \VUgras{نماذج} جميلة معلَّقة على الجدار
و\VUgras{مصباح} \textsuperscript{(62)} \VUgras{بغطاء}، متدلٍّ من السقف.
و\VUgras{المعلم} \textsuperscript{(60)} صبور جدًا؛ يصحح \VUgras{رسوم}
\textsuperscript{(61)} التلاميذ ويعلّمهم رسم \VUgras{تمثال نصفي} \textsuperscript{(59)}
عن الطبيعة.

%%K t01-16-1 p
16. --- و\VUgras{فصل الموسيقى} يبدو ممتعًا جدًا. فيه يُتعلَّم قراءة
\VUgras{النوتة}، وغناء \VUgras{درجات السلَّم} \textsuperscript{(67)} المكتوبة على
\VUgras{المُدرَّج} (دو، ري، مي، فا، صول، لا، سي\,=\,C. D. E. F. G. A. B.)،
ومعرفة \VUgras{المفاتيح}، وغناء \VUgras{ألحان} بديعة. وتوضع النوتات على
\VUgras{الحامل} \textsuperscript{(65)}. وأحد التلاميذ \VUgras{يعزف على البيانو}
\textsuperscript{(64)}، و\VUgras{معلم الموسيقى} \textsuperscript{(66)}
\VUgras{يضبط الإيقاع} \textsuperscript{(68)}.

%%K t01-17-1 p
17. --- ونلمح ركنًا من ورشة \VUgras{الأشغال اليدوية} \textsuperscript{(69)}،
حيث يستطيع الأطفال أن يتعلموا مسح الخشب وبرد الحديد. وليست في كل مدرسة
ورشة كهذه.

%%K t01-tit-7 sub
\VUsaut{5.0mm}
\VUpk{10.2pt}{40}{فصل العلوم.}
\VUsaut{3.6mm}

%%K t01-18-1 p
18. --- معلم الرياضيات يعلّم التلاميذ \VUgras{الهندسة}
و\VUgras{الحساب والعمليات} و\VUgras{النظام المتري}. وعلى اللوحة نرى الأشكال
الهندسية الرئيسية: \VUgras{دائرة} \textsuperscript{(a)} و\VUgras{مربعًا}
\textsuperscript{(b)} و\VUgras{مثلثًا} \textsuperscript{(c)} و\VUgras{خطًا}
\textsuperscript{(d)}.

%%K t01-18-2 p
ونرى أيضًا \VUgras{عملية}؛ وهي ليست \VUgras{جمعًا} ولا \VUgras{طرحًا} ولا
\VUgras{ضربًا}: إنها \VUgras{قسمة} \textsuperscript{(e)}.

%%K t01-19-1 p
19. --- وعلى لوحة النظام المتري نرى: \VUgras{مترًا} \textsuperscript{(a)}
و\VUgras{ميزانًا} \textsuperscript{(b)} و\VUgras{صنجاته}، و\VUgras{لترًا}
\textsuperscript{(e)} و\VUgras{ديكالترًا} \textsuperscript{(f)} و\VUgras{ديسيلترًا}
و\VUgras{مترًا مكعبًا} \textsuperscript{(d)}.

%%K t01-19-2 p
ويدرس التلاميذ أيضًا \VUgras{العلوم الطبيعية} \textsuperscript{(11)} --- علم
الحيوان \textsuperscript{(a)} وعلم النبات \textsuperscript{(b)} وعلم الأرض
\textsuperscript{(c)} --- و\VUgras{التاريخ} \textsuperscript{(12)}.

%%K t01-19-3 p
وفي الخزانة أدوات \VUgras{الفيزياء} \textsuperscript{(15)} و\VUgras{الكيمياء}
\textsuperscript{(16)}.

%%K t01-19-4 p
وفوق الخزانة: \VUgras{كرة} \textsuperscript{(14)} و\VUgras{مخروط}
و\VUgras{مُجسَّم أرضي} \textsuperscript{(13)}.

%%K t01-19-5 p
وفوق اللوحات \VUgras{خريطة} تمثّل أجزاء العالم الخمسة: \VUgras{أوروبا}
\textsuperscript{(g)} و\VUgras{آسيا} \textsuperscript{(e)} و\VUgras{إفريقيا}
\textsuperscript{(f)} و\VUgras{أمريكا} \textsuperscript{(ab)} و\VUgras{أوقيانوسيا}
\textsuperscript{(h)}، والمحيطين العظيمين: \VUgras{الهادئ} \textsuperscript{(c)}
و\VUgras{الأطلسي} \textsuperscript{(d)}.
\VUsaut{6.0mm}
\VUfilet{22mm}
